\documentclass[11pt,a4paper,oneside]{article}

% --- Core Packages ---
\usepackage[utf8]{inputenc}
\usepackage[T1]{fontenc}
\usepackage[english]{babel}
\usepackage{amsmath, amssymb, amsfonts}
\usepackage{graphicx}
\usepackage{geometry}
\usepackage{xcolor}
\usepackage{hyperref}
\usepackage{booktabs}
\usepackage{caption}
\usepackage{enumitem}

% --- Document Layout ---
\geometry{
    a4paper,
    total={170mm,257mm},
    left=20mm,
    top=20mm,
}

% --- Custom Colors & Links ---
\definecolor{uetgreen}{HTML}{2D5A27}
\hypersetup{
    colorlinks=true,
    linkcolor=uetgreen,
    filecolor=magenta,
    urlcolor=uetgreen,
    citecolor=uetgreen,
}

% --- Title Information ---
\title{\textbf{Universal Equilibrium Theory: [Topic Title]}}
\author{Unity Equilibrium Group \\ \small \href{https://github.com/unityequilibrium}{github.com/unityequilibrium}}
\date{\today}

\begin{document}

\maketitle

\begin{abstract}
\textbf{Problem:} [State the bounded scientific or technical question.] \\
\textbf{Method:} [State the model, derivation, or benchmark method used.] \\
\textbf{Result:} [State the main measurable outcome conservatively.] \\
\textbf{Scope:} [State whether the evidence is internal benchmark, external replication, or another bounded class.]
\end{abstract}

\section{Introduction}
State the specific problem, why it matters, and how this paper is scoped. Keep the framing
focused on one topic-scale question rather than the whole repository ambition.

\section{Theoretical Framework}
Present the exact model, derivation, or conceptual framework used for this topic.
\begin{equation}
    \Omega(C) = \alpha(C^2 - C0^2)^2 + \beta C + \kappa (\nabla C)^2
\end{equation}
State assumptions explicitly and avoid stronger wording than the evidence supports.

\subsection{Conceptual Logic Map}
Use a figure only if it clarifies the method or evidence flow.

\begin{figure}[h]
    \centering
    % \includegraphics[width=0.8\textwidth]{../Result/01_Showcase/Social_Proof.png}
    \fbox{
        \begin{minipage}{0.8\textwidth}
            \centering
            \vspace{2cm}
            [INSERT METHOD OR EVIDENCE FLOW DIAGRAM HERE] \\
            \textit{(Problem $\rightarrow$ Method $\rightarrow$ Evidence)}
            \vspace{2cm}
        \end{minipage}
    }
    \caption{Conceptual architecture for [Topic Area].}
\end{figure}

\section{Data Integrity and Sources}
\textbf{Primary data source:} [Name]. \\
\textbf{Reference ID:} [DOI / Publication URL]. \\
\textbf{Preprocessing note:} [How the local working copy relates to the upstream source.]

\section{Methodology}
Describe the workflow so another reviewer can understand the code, data, metrics, and
thresholds used.
\begin{itemize}
    \item \textbf{Engine or method path:} [Engine Name or script]
    \item \textbf{Verification path:} [Proof or benchmark script]
    \item \textbf{Scale or regime:} [Planck / Macroscopic / Astrophysical / Other]
    \item \textbf{Fitting or calibration used:} [Yes/No, with note]
\end{itemize}

\section{Results}
Present results with explicit metrics, baselines, and uncertainty-aware interpretation.
\begin{table}[h!]
  \centering
  \caption{Comparison between the studied approach and baseline.}
  \begin{tabular}{@{}lll@{}}
    \toprule
    Metric & Baseline & This work \\ \midrule
    Primary metric & [Value] & [Value] \\
    Threshold status & [Value] & [Pass/Fail with condition] \\
    Notes & [Context] & [Context] \\ \bottomrule
  \end{tabular}
\end{table}

\section{Discussion}
Explain what the result means, what alternative explanations remain, and what still limits
interpretation.

\section{Limitations}
State what is still speculative, internally bounded, fitted, or not yet externally
replicated.

\section{Conclusion}
Summarize the bounded contribution and avoid claiming more than the paper actually supports.

\section*{Open Source and Reproducibility}
The code, data manifests, and verification scripts are available at: \\
\url{https://github.com/unityequilibrium/Lab_uet_harness}

\section*{Reproducibility Notes}
\begin{itemize}
    \item Repository version: [version]
    \item Verification command: [command]
    \item Artifact path: [Result/artifacts/...json]
\end{itemize}

\bibliographystyle{plain}
\section*{Key Academic References}
\begin{enumerate}
    \item [Author Name], \textit{[Title]}, DOI: [10.xxxx/xxx]
    \item [Benchmark or comparison source]
\end{enumerate}

\end{document}
